% Pursuing Stacks (cotes 134-2 to 134-8): archive cote 134-7, site page 55 (PDF page 56), typescript page 497.
% typescript. Draft compared: ps6.tex:2719-2783.
% AI-transcribed by Claude Opus 5.5 (claude-opus-5-5, adaptive thinking, effort high), Message Batches API, batch ps4 (msgbatch_012Aoc1NuskATPsyBeKSYK7J), 2026-11-12; not yet checked by eye.
\scanpage{134-7/55}
\mspage{497}

I do'nt have,  I must confess,  any direct description of \inserted{such} an isomorphism
(9),  valid for any semisimplicial bundle $X_{*}$  say satisfying the assumti-%
ons
\[ X_0=X_1= e \ , \quad X_n \text{ smooth for any n } \hand{\unsure{,}}\text{(maybe flat is enough) ,} \]
plus possibly (if needed) a Kan type condition.  However,  we have
such isomorphisms (9) in a tautological way,  when $X_{*}$  comes \xout{\unsure{fxxix}} in the
usual way from a chain complex in Ab$_k$  with projective components,  hence
also when $X_{*}$  admits a Postnikov-type dévissage into ``abelian'' pieces as
above.  If we admit that any $X_{*}$  satisfying the assumptions is homotopic
to one admitting such a dévissage,  the isomorphisms (9) should follow,
except of course that extra work would be still needed to get naturality
of (9).

I have the feeling however that,  besides the specific abelianiza-%
tion functor in the schematic context,  \uline{formula} (9) \uline{should be} \xout{viewed as} \inserted[typed]{made a cor-}
\xout{beingxx``the''xnaturalx\uline{definition}} \inserted[typed]{nerstone of a theory of schematic}  \uline{homotopy types},  \uline{and serve as}``\uline{the}''
\uline{natural definition of the homotopy invariants} of a model $X_{*}$ ,  within the
context of schematic models and without any need a priori to \xout{subordinatex}
\xout{these} tie them up with,  let alone subordinate their study\hand{\unsure{,}} to,  \xout{corres-}
\xout{pondingxdiscrxhomotopy}invariants of  the corresponding discrete homotopy
type $X_{*}(k)$.  Accordingly,  \xout{quasi} weak equivalences should be defined \xout{as}
\xout{maps} (for semisimplicial bundles satisfying the suitable assumptions at
any rate)  as maps inducing isomorphisms for the $\hand{\pi}_i$ invariants,  namely
inducing quasi-isomorphisms for the corresponding Lie chain complexes.
It is immediately checked that this implies that the corresponding map
for ``formal homology'' namely
\[ L_{pt}(X_{*}) \hand{\to} L_{pt}(Y_{*}) \]
is then a quasi-isomorphism too,  when viewed as a map of chain complexes,
and hopefully the same should hold for ``schematic homology''
\[ L(X_{*}) \hand{\to} L(Y_{*}) \ \ , \]
and of course one would expect the converse statements to hold too.

I would like to comment a little on the significance of formula (9).
\xout{To my k} As far as I know,  this is the only fairly general formula,  not
reducing to an ``abelian'' case,  where the homotopy groups $\hand{\pi}_i$  appear
as just the $H_i$  invariants of a suitable chain complex,  def\hand{i}ned up to
unique isomorphism in the relevant derived category.  This chain complex
comes here,  moreover,  with an amazingly simple description,  of immedia-%
te geometrical significance,  and suggestive of \xout{very p} relationships with

\margin{55}
